%% ma: L12-B11-0004
%% lop: 12
%% chuong: 4
%% bai: 11
%% dang: 12.11.4
%% loai: DS
%% muc: [NB, TH, VD, VD]
%% dap_an: "ĐSĐĐ"
%% nguon: "(NBV)-ĐỀ SỐ 5-MỖI NGÀY 1 ĐỀ THI 2025.pdf (D05, MỖI NGÀY 1 ĐỀ - Đề số 5), Phần II Câu 4"
%% kien_thuc: "Bài 11 (nguyên hàm của hàm số mũ, hằng số C từ điều kiện ban đầu); Bài 20 lớp 11 (hàm số mũ)"
%% trang_thai: cho-duyet
%% ly_do: "Dạng toán mới đề xuất 12.11.4: tìm S(t) từ tốc độ tăng dân số S'(t) = k·e^{rt} (nguyên hàm, Đúng/Sai)"
%% ghi_chu: "Bỏ ảnh minh họa đám đông: ảnh trang trí, không chứa dữ kiện toán"
\begin{ex}
Vào năm $2014$, dân số nước ta khoảng $90{,}7$ triệu người. Giả sử, dân số nước ta sau $t$ năm được xác định bởi hàm số $S(t)$ (đơn vị: triệu người), trong đó tốc độ gia tăng dân số được cho bởi $S'(t)=1{,}2698e^{0{,}014t}$, với $t$ là số năm kể từ năm $2014$, $S'(t)$ tính bằng triệu người/năm.
\choiceTF
{\True $S(t)$ là một nguyên hàm của $S'(t)$.}
{$S(t)=90{,}7e^{0{,}014t}+90{,}7$.}
{\True Theo công thức trên, tốc độ tăng dân số nước ta năm $2034$ (làm tròn đến hàng phần mười của triệu người/năm) khoảng $1{,}7$ triệu người/năm.}
{\True Theo công thức trên, dân số nước ta năm $2034$ (làm tròn đến hàng đơn vị của triệu người) là khoảng $120$ triệu người.}
\loigiai{a) $S'(t)$ là đạo hàm của $S(t)$ nên $S(t)$ là một nguyên hàm của $S'(t)$. Đúng.
b) $S(t)=\displaystyle\int1{,}2698e^{0{,}014t}\,\mathrm dt=\dfrac{1{,}2698}{0{,}014}e^{0{,}014t}+C=90{,}7e^{0{,}014t}+C$. Vì $S(0)=90{,}7$ nên $C=0$, tức $S(t)=90{,}7e^{0{,}014t}$. Công thức ở ý b) cho $S(0)=181{,}4\ne90{,}7$. Sai.
c) Năm $2034$ ứng với $t=20$: $S'(20)=1{,}2698e^{0{,}28}\approx1{,}68\approx1{,}7$ (triệu người/năm). Đúng.
d) $S(20)=90{,}7e^{0{,}28}\approx120{,}0\approx120$ (triệu người). Đúng.}
\end{ex}
